% ===============================
% proposal_content.tex
% เนื้อหาเสนอเค้าโครงโครงงาน
% แก้ไขเนื้อหาในไฟล์นี้
% ===============================

% -------------------------------------------------------
\section{ชื่อหัวข้อโครงงาน}
% -------------------------------------------------------

ภาษาไทย การพัฒนาโมเดลจำแนกประเภทข้อความร้องเรียนของลูกค้า

ภาษาอังกฤษ Customer Complaint Text Classification

% -------------------------------------------------------
\section{ที่มาและความสำคัญของปัญหา}
% -------------------------------------------------------

ปัจจุบันธุรกิจอีคอมเมิร์ซ (E-commerce) ในประเทศไทยมีการเติบโตอย่างต่อเนื่อง สอดคล้องกับรายงานของ Statista ที่ระบุว่ามูลค่าตลาดอีคอมเมิร์ซของไทยมีแนวโน้มเพิ่มสูงขึ้นในทุกปี~\cite{statista2024ecommerce} การขยายตัวดังกล่าวส่งผลให้ปริมาณการติดต่อและข้อร้องเรียนจากลูกค้าเพิ่มสูงขึ้นตามไปด้วย~\cite{roumeliotis2025think} โดยข้อร้องเรียนเหล่านี้มีความหลากหลายและครอบคลุมหลายประเด็น เช่น การจัดส่งสินค้าผิดพลาด สินค้าชำรุด สินค้าไม่ครบจำนวน การขอใบกำกับภาษี การขอคืนเงินหรือคืนสินค้า และการติดตามสถานะการจัดส่ง เป็นต้น ซึ่งหากองค์กรต้องอาศัยพนักงาน (Human Agent) ในการอ่านและคัดแยกหมวดหมู่ข้อความด้วยตนเองทีละรายการ ย่อมก่อให้เกิดความล่าช้าในการตอบสนองปัญหา และอาจส่งผลกระทบเชิงลบต่อประสบการณ์และความพึงพอใจของลูกค้าโดยตรง

อย่างไรก็ตาม แนวทางที่ใช้ในปัจจุบันสำหรับการคัดแยกข้อความร้องเรียนแบบอัตโนมัติยังมีข้อจำกัดหลายประการ ได้แก่ (1)~ระบบที่อาศัยคำสำคัญ (Keyword-based System) ซึ่งต้องอาศัยผู้เชี่ยวชาญในการกำหนดกฎ แม้จะมีความเร็วในการประมวลผลสูง แต่ขาดความยืดหยุ่นเมื่อเผชิญกับภาษาพูด คำสแลง หรือประโยคที่มีความซับซ้อนกำกวม (2)~โมเดลการเรียนรู้ของเครื่องแบบดั้งเดิม (Traditional Machine Learning) เช่น Naive Bayes หรือ SVM ที่จำเป็นต้องอาศัยชุดข้อมูลสำหรับฝึกสอน (Labeled Data) ปริมาณมหาศาล อีกทั้งยังมีข้อจำกัดด้านประสิทธิภาพเมื่อนำมาประมวลผลภาษาไทย และ (3)~ระบบการจัดการตั๋วปัญหา (Ticketing System) ที่แม้จะมีการจัดหมวดหมู่ที่เป็นระเบียบ แต่ก็ยังคงต้องอาศัยมนุษย์ในการคัดกรองเบื้องต้นในหลายกรณี

จากข้อจำกัดดังกล่าว โครงงานนี้จึงเล็งเห็นถึงความท้าทายในการพัฒนาโมเดลจำแนกประเภทข้อความภาษาไทย เนื่องจากภาษาไทยมีทรัพยากรข้อมูลที่จำกัด ส่งผลให้โมเดลทั่วไปมีประสิทธิภาพลดลง~\cite{sriphon2025} ประกอบกับการขาดแคลนชุดข้อมูลการร้องเรียนของลูกค้าภาษาไทยที่ถูกจัดหมวดหมู่ไว้อย่างสมบูรณ์ โครงงานนี้จึงมุ่งศึกษาและประยุกต์ใช้โมเดลภาษาขนาดใหญ่ (Large Language Models: LLMs) ด้วยเทคนิค Zero-shot Classification ซึ่งสามารถจำแนกหมวดหมู่ข้อความได้โดยไม่จำเป็นต้องใช้ชุดข้อมูลฝึกสอนจำนวนมาก และนำมาเปรียบเทียบกับโมเดลแบบดั้งเดิมที่ผ่านการฝึกสอน (Supervised Models) เพื่อค้นหาแนวทางที่มีประสิทธิภาพและเหมาะสมที่สุดสำหรับการจัดการข้อความร้องเรียนภาษาไทย


% -------------------------------------------------------
\section{วัตถุประสงค์ของโครงงาน}
% -------------------------------------------------------

\begin{enumerate}
    \item เพื่อพัฒนาระบบจำแนกข้อความร้องเรียนของลูกค้า
    \item เพื่อศึกษาประสิทธิภาพของ Large Language Models ในการทำ Zero-shot Classification
    \item เพื่อเปรียบเทียบประสิทธิภาพระหว่างโมเดล LLM แบบ Zero-shot Classification กับโมเดลที่ผ่านการฝึกจากชุดข้อมูล
\end{enumerate}

% -------------------------------------------------------
\section{ทฤษฎีและงานวิจัยที่เกี่ยวข้อง}
% -------------------------------------------------------

\subsection{ทฤษฎีที่เกี่ยวข้อง}

\subsubsection{การประมวลผลภาษาธรรมชาติและการตัดคำภาษาไทย (Thai Word Tokenization)}
การตัดคำ (Word Tokenization) เป็นขั้นตอนสำคัญในการประมวลผลภาษาธรรมชาติ โดยเฉพาะภาษาไทยที่ไม่มีการเว้นวรรคระหว่างคำอย่างชัดเจน การตัดคำที่นิยมใช้มี 2 แนวทางหลัก ได้แก่ การตัดคำแบบอิงพจนานุกรม (Dictionary-based) เช่น อัลกอริทึม Maximal Matching ที่ใช้ในไลบรารี PyThaiNLP~\cite{phatthiyaphaibun2023pythainlp} ซึ่งให้ผลลัพธ์ที่อ่านเข้าใจง่ายแต่อาจมีปัญหาเมื่อพบคำศัพท์ใหม่ (Out-of-Vocabulary: OOV) และแนวทางการตัดคำระดับหน่วยย่อย (Subword Tokenization) เช่น SentencePiece ซึ่งใช้หลักการทางสถิติในการเรียนรู้หน่วยย่อยของคำ ทำให้สามารถจัดการกับคำศัพท์ใหม่ได้ดีกว่า ซึ่งเป็นแนวทางหลักที่นิยมใช้ร่วมกับโมเดลภาษาขนาดใหญ่ในปัจจุบัน

\subsubsection{โมเดลการเรียนรู้ของเครื่องสำหรับการจำแนกประเภทข้อความ}
สถาปัตยกรรม Transformer ถือเป็นรากฐานสำคัญของโมเดลที่ใช้ในโครงงานนี้ โดยแบ่งเป็น 2 กลุ่มหลัก ได้แก่:
\begin{enumerate}[label=(\arabic*)]
    \item โมเดลภาษาเฉพาะภาษาไทย (Thai-centric Language Models): ได้แก่ WangchanBERTa~\cite{lowphansirikul2021wangchanberta} ซึ่งพัฒนาบนสถาปัตยกรรม RoBERTa และฝึกสอนด้วยชุดข้อมูลภาษาไทยขนาดใหญ่กว่า 78.5 GB โดยใช้ SentencePiece ทำให้มีความเข้าใจบริบทภาษาไทยได้อย่างมีประสิทธิภาพ นอกจากนี้ยังมีโมเดลขนาดใหญ่อย่าง Typhoon2-7B~\cite{sriphon2025} ที่ออกแบบให้รองรับการสั่งงานด้วยภาษาไทย
    \item โมเดลภาษารองรับหลายภาษา (Multilingual Models): ได้แก่ Multilingual BERT (mBERT) และ LLM เชิงพาณิชย์ เช่น GPT-4o และ Claude 3.5 Sonnet ซึ่งมีความสามารถในการแปลและทำความเข้าใจได้หลายภาษาแบบ Zero-shot 
\end{enumerate}

\subsubsection{ตัวชี้วัดประสิทธิภาพ (Evaluation Metrics)}
การประเมินประสิทธิภาพของโมเดลจำแนกประเภทข้อความ โดยเฉพาะข้อมูลที่มีความไม่สมดุล (Imbalanced Data) มักใช้ตัวชี้วัดดังนี้:
\begin{enumerate}[label=(\arabic*)]
    \item Accuracy (ความแม่นยำ): สัดส่วนที่โมเดลสามารถทำนายประเภทได้ถูกต้องจากข้อมูลทั้งหมด
    \item Macro F1-Score: ค่าเฉลี่ยของ F1-Score จากทุกคลาส โดยให้น้ำหนักแต่ละคลาสเท่า ๆ กัน ช่วยให้เห็นประสิทธิภาพรวมได้ดีในกรณีที่บางคลาสมีตัวอย่างน้อย
    \item Confusion Matrix: ตารางแสดงการเปรียบเทียบระหว่างประเภทที่โมเดลทำนายกับประเภทที่เป็นจริง ช่วยวิเคราะห์ว่าโมเดลมักสับสนระหว่างคลาสใด (Misclassification)
\end{enumerate}

\subsection{งานวิจัยที่เกี่ยวข้อง}

\subsubsection{Text Classification in the LLM Era — Where do we stand? \cite{vajjala2024}}

งานวิจัยของ Vajjala และ Shimangaud \cite{vajjala2024} ได้เปรียบเทียบวิธีการจำแนกข้อความ 4 รูปแบบ ได้แก่ Zero-shot, Few-shot Fine-tuning, Synthetic Data-based Classification และ Fully Supervised Classification บนชุดข้อมูล 32 ชุด ครอบคลุม 8 ภาษา ผลการทดลองพบว่า Zero-shot ให้ผลดีในงาน Sentiment Analysis แต่สำหรับงานจำแนกประเภทอื่น Few-shot Fine-tuning และ Fully Supervised ยังให้ความแม่นยำสูงกว่า และประสิทธิภาพของ LLM แตกต่างกันไปตามภาษา งานวิจัยนี้ไม่ได้มุ่งเน้นข้อมูลร้องเรียนของลูกค้าโดยเฉพาะและไม่ได้ศึกษาการออกแบบ Prompt อย่างละเอียด แต่ให้ภาพรวมสำคัญของแนวทางที่นำมาใช้ในโครงงานนี้

\subsubsection{Adaptable and Reliable Text Classification using Large Language Models \cite{wang2023adaptable}}

งานวิจัยของ Wang และคณะ \cite{wang2023adaptable} เปรียบเทียบประสิทธิภาพระหว่าง Machine Learning, Deep Learning และ LLM ภายในกรอบเดียวกัน โดยทดลอง 3 รูปแบบ คือ Zero-shot, Few-shot และ Fine-tuning พบว่า LLM จำแนกข้อความได้ดีแม้ไม่มีข้อมูลฝึกเพิ่มเติม และการออกแบบ Prompt มีผลต่อประสิทธิภาพอย่างชัดเจน โดยเฉพาะกับชุดข้อมูลขนาดเล็ก LLM มีความยืดหยุ่นสูงกว่าโมเดลดั้งเดิม ข้อจำกัดคือยังไม่ได้ทดลองกับข้อมูลภาษาไทย และงานจำแนกข้อความร้องเรียนของลูกค้าโดยตรง แนวคิดเรื่องคุณภาพของ Prompt จากงานวิจัยนี้จึงถูกนำมาใช้เป็นแนวทางออกแบบ Prompt ของโครงงาน

\subsubsection{Think Before You Classify: The Rise of Reasoning LLMs for Consumer Complaint Classification \cite{roumeliotis2025think}}

งานวิจัยของ Roumeliotis และคณะ \cite{roumeliotis2025think} เปรียบเทียบโมเดล GPT-4o, Claude, Gemini, DeepSeek และ Reasoning LLM ในการจำแนกข้อความร้องเรียนจากฐานข้อมูล Consumer Complaint Database ด้วยวิธี Zero-shot Prompting ภายใต้ Prompt เดียวกัน ผลการทดลองพบว่าโมเดลที่มีความสามารถด้าน Reasoning ให้ผลดีกว่า LLM ทั่วไป และ Claude มีประสิทธิภาพสูงสุด

\subsubsection{DeepSeek and GPT Fall Behind: Claude Leads in Zero-Shot Consumer Complaints Classification \cite{roumeliotis2025deepseek}}

Roumeliotis และคณะ \cite{roumeliotis2025deepseek} ทดสอบเปรียบเทียบ Claude, GPT, DeepSeek และ Gemini บนชุดข้อมูล Consumer Complaint Database ยืนยันว่า Claude มีประสิทธิภาพสูงสุดในการจำแนกข้อความร้องเรียนแบบ Zero-shot และคุณภาพของ Prompt ส่งผลต่อประสิทธิภาพของทุกโมเดลอย่างชัดเจน

\subsubsection{Evaluating the impact of model scale and prompting strategies on corruption allegation classification using Thai-specialized Typhoon2 language models \cite{sriphon2025}}

Sriphon และคณะ \cite{sriphon2025} ศึกษาผลของขนาดโมเดลและรูปแบบ Prompt (zero-shot, one-shot, two-shot) ของโมเดลภาษาไทย Typhoon2 ต่อการจำแนกข้อร้องเรียนเรื่องทุจริตที่ยื่นต่อสำนักงานคณะกรรมการป้องกันและปราบปรามการทุจริตแห่งชาติ (ป.ป.ช.) ผลการทดลองพบว่าโมเดลขนาดใหญ่กว่า (Typhoon2-7B) ร่วมกับ Prompt แบบ two-shot ให้ผลลัพธ์ที่ดีที่สุดในกลุ่ม LLM แต่โมเดล Machine Learning แบบดั้งเดิม (Random Forest, XGBoost) ยังคงให้ผลลัพธ์ที่ดีกว่าในภาพรวม

\subsubsection{Ensembling Prompting Strategies for Zero-shot Hierarchical Text Classification with LLMs \cite{xia2024ensembling}}

Xia และคณะ \cite{xia2024ensembling} เสนอแนวทาง Prompt Ensembling คือ การรวมผลลัพธ์จาก Prompt หลายรูปแบบ (Direct, Instruction, Reasoning) เพื่อเพิ่มความแม่นยำของการจำแนกข้อความแบบ Zero-shot โดยไม่ต้องฝึกโมเดลเพิ่มเติม ผลการทดลองพบว่าวิธีนี้ช่วยลดความผิดพลาดในข้อความที่มีความกำกวมหรือซับซ้อน ข้อจำกัดคือใช้เวลาประมวลผลเพิ่มขึ้นและยังไม่ได้ทดลองกับภาษาไทย

\subsubsection{Large language models are zero-shot text classifiers \cite{wang2023zeroshot}}

Wang และคณะ \cite{wang2023zeroshot} แสดงให้เห็นว่า LLM สามารถทำหน้าที่เป็น Zero-shot Text Classifier ได้อย่างมีประสิทธิภาพในหลายงาน เช่น Sentiment Analysis, Topic Classification และ Intent Classification โดยไม่ต้องฝึกโมเดลเพิ่มเติม แม้ Fine-tuning จะให้ผลแม่นยำกว่าในบางงาน แต่ Zero-shot มีข้อดีเปรียบด้านความสะดวกและความรวดเร็ว

% -------------------------------------------------------
\section{วิธีดำเนินการวิจัย}
% -------------------------------------------------------

\subsection{การเก็บรวบรวมข้อมูล}
    
รวบรวมข้อความร้องเรียนของลูกค้าจากระบบบริการลูกค้าของร้านค้าออนไลน์ (ช่องทาง Live Chat และ Inbox) ในช่วงกรอบเวลา 3 เดือน (พฤษภาคม - กรกฎาคม 2569) โดยตั้งเป้าหมายจำนวนตัวอย่างไว้ที่ 1,000 ข้อความ ที่ครอบคลุม 11 ประเภทตาม label ที่กำหนด ข้อมูลจะผ่านการคัดกรองเบื้องต้นเพื่อตัดข้อความสแปมหรือข้อความสั้นที่ไม่มีใจความสำคัญออก (เกณฑ์ความน่าเชื่อถือเบื้องต้น) และดำเนินการ anonymize ลบข้อมูลส่วนบุคคล (PDPA) เพื่อปกป้องความเป็นส่วนตัวของลูกค้า

\subsection{การเตรียมข้อมูล (Data Preprocessing และ Annotation)}
    
ทำความสะอาดข้อความ เช่น ลบ Emoji, URL, และอักขระพิเศษที่ไม่เกี่ยวข้อง จากนั้นดำเนินการ annotation โดยผู้ติด label อย่างน้อย 2 คนต่อ 1 ข้อความ และวัด Inter-annotator Agreement ด้วย Cohen's Kappa เพื่อยืนยันความสม่ำเสมอ จากนั้นแบ่งชุดข้อมูลเป็น Train 70\% / Validation 15\% / Test 15\% โดยรักษาสัดส่วน label ให้สมดุล (Stratified Split)

\subsection{การพัฒนาโมเดลจำแนกข้อความ (Supervised Baseline)}
    
ฝึกโมเดล baseline ด้วยชุดข้อมูล Train โดยใช้ WangchanBERTa~\cite{sriphon2025} ซึ่งเป็นโมเดลภาษาไทย และ mBERT (Multilingual BERT) เป็นตัวเปรียบเทียบ โดยใช้ library HuggingFace Transformers และ fine-tune ด้วย GPU บน Google Colab

\subsection{การทดสอบด้วย Zero-shot Classification ผ่าน LLM}
    
ออกแบบ Prompt ใน 3 รูปแบบ ได้แก่ Direct Prompt, Instruction Prompt และ Reasoning Prompt~\cite{xia2024ensembling} และทดสอบกับโมเดล LLM ที่เลือกใช้ ได้แก่ Claude 3.5 Sonnet (Anthropic API), GPT-4o (OpenAI API) และ Typhoon2-7B~\cite{sriphon2025} (สำหรับภาษาไทยโดยเฉพาะ) โดยส่ง prompt ผ่าน API และบันทึกผลลัพธ์บน Test set ชุดเดียวกันกับ supervised baseline

\subsection{การประเมินผลและเปรียบเทียบ}
    
วัดประสิทธิภาพด้วย Accuracy และ Macro F1-score บน Test set วิเคราะห์ Confusion Matrix รายประเภท เพื่อระบุ label ที่โมเดลจำแนกได้ยาก และเปรียบเทียบประสิทธิภาพระหว่างโมเดล supervised (WangchanBERTa, mBERT) กับ Zero-shot LLM (Claude, GPT-4o, Typhoon2) รวมถึงวิเคราะห์ผลกระทบของรูปแบบ Prompt ต่อประสิทธิภาพ


% -------------------------------------------------------
\section{การวิเคราะห์และออกแบบระบบ}
% -------------------------------------------------------

เนื่องจากโครงงานนี้มุ่งเน้นการวิจัยและพัฒนาโมเดลปัญญาประดิษฐ์ (AI/ML Research) ระบบที่จะพัฒนาจึงอยู่ในรูปแบบของ Data Processing Pipeline สำหรับการจำแนกข้อความอัตโนมัติ โดยมีรายละเอียดการวิเคราะห์และออกแบบสถาปัตยกรรมดังนี้

\subsection{ผู้ใช้งานระบบ (Actors) และสิทธิ์การใช้งาน}

\subsubsection{พนักงานฝ่ายบริการลูกค้า (CS Staff)}
เป็นผู้ส่งต่อข้อมูลข้อความร้องเรียนดิบที่ส่งเข้ามาในแต่ละวันเข้าสู่ระบบ และรับผลลัพธ์การจำแนกประเภทเพื่อนำไปแจกจ่ายงาน (Ticket Routing) ให้แผนกที่เกี่ยวข้องต่อไป (Use case: นำเข้าข้อมูล, ดูผลลัพธ์การจำแนก)

\subsubsection{ผู้พัฒนาระบบ (AI Engineer/Data Scientist)}
เป็นผู้ดูแลระบบ อัปเดตชุดข้อมูล ทดสอบ prompt และบำรุงรักษาโมเดลให้มีความแม่นยำ (Use case: ปรับปรุงโมเดล, ประเมินประสิทธิภาพ)

\subsection{สถาปัตยกรรมและกระบวนการทำงาน (Pipeline Architecture)}
การทำงานของระบบประมวลผล (Data Flow) ถูกออกแบบเป็นไปป์ไลน์ 4 ขั้นตอนหลักดังนี้:

\subsubsection{Data Ingestion (การรับข้อมูล)}
ระบบรับข้อมูลข้อความร้องเรียนจากฐานข้อมูลลูกค้าในรูปแบบไฟล์ (เช่น CSV) หรือดึงข้อมูลผ่าน API จากระบบ Ticketing เดิม

\subsubsection{Data Preprocessing (การเตรียมข้อมูล)}
โมดูลทำความสะอาดข้อความ ลบคำหยาบ และกระบวนการ Anonymization/De-identification เพื่อปกปิดข้อมูลส่วนบุคคล (PII) ของลูกค้าก่อนส่งให้โมเดลประมวลผล

\subsubsection{Inference Engine (เอนจินจำแนกข้อความ)}
ข้อมูลที่พร้อมใช้งานจะถูกส่งเข้าสู่แกนหลักของระบบเพื่อทำนายประเภท
\begin{enumerate}[label=(\arabic*)]
    \item \textit{กรณีใช้ Supervised Model}: ข้อความถูกส่งเข้าโมเดล WangchanBERTa ที่ผ่านการ Fine-tune และ Deploy ไว้บนเซิร์ฟเวอร์
    \item \textit{กรณีใช้ Zero-shot LLM}: ข้อความถูกประกอบเข้ากับ Prompt ที่ออกแบบไว้ และส่งไปยัง API ของ Large Language Model (เช่น Claude 3.5 Sonnet หรือ GPT-4o)
\end{enumerate}

\subsubsection{Post-processing \& Output}
ผลลัพธ์ที่ได้จากการทำนาย (Predicted Label) จะถูกแนบกลับไปกับข้อความต้นฉบับ และส่งออกเป็นรายงานสรุป (Report) หรือส่งคืนกลับไปยังระบบ Ticketing เพื่อให้เจ้าหน้าที่ CS นำไปดำเนินการแก้ไขปัญหาให้ลูกค้าต่อไป

% -------------------------------------------------------
\section{ขอบเขตของโครงงาน}
% -------------------------------------------------------

\subsection{ข้อมูล}

\noindent\textbf{ตารางที่~1} ตารางข้อมูล Label

\vspace{0.2cm}
\noindent
\setlength{\LTleft}{0pt}
\begin{longtable}{|p{3.4cm}|p{3.1cm}|p{7.7cm}|}

% --- หัวตารางหน้าแรก (ไม่มี caption ซ้ำ) ---
\hline
\textbf{Label} & \textbf{ชื่อไทย} & \textbf{ตัวอย่างข้อความจากลูกค้า} \\
\hline
\endfirsthead

% --- หัวตารางหน้าถัดไป ---
\multicolumn{3}{@{}l@{}}{\textbf{ตารางที่~1} ตารางข้อมูล Label (ต่อ)} \\[2pt]
\hline
\textbf{Label} & \textbf{ชื่อไทย} & \textbf{ตัวอย่างข้อความจากลูกค้า} \\
\hline
\endhead

% --- เส้นปิดท้าย (ทุกหน้า) ---
\hline
\endfoot

% --- ข้อมูล ---
Wrong Item              & ได้รับสินค้าผิด   & สั่งเสื้อสีดำแต่ได้รับสีขาวมาแทนคะ \\
\hline
Damaged Item            & สินค้าชำรุด       & กล่องมาถึงบุบและของข้างในแตกคะ \\
\hline
Missing Item            & ของไม่ครบ         & สั่งไป 3 ชิ้น แต่ได้รับมาแค่ 2 ชิ้นครับ \\
\hline
Invoice Request         & ขอใบกำกับภาษี     & รบกวนขอใบกำกับภาษีเต็มรูปแบบสำหรับออเดอร์นี้ด้วยคะ \\
\hline
Return                  & ขอคืนสินค้า       & อยากคืนสินค้าเพราะไม่ตรงกับที่สั่ง ต้องทำอย่างไรบ้างคะ \\
\hline
Exchange                & ขอเปลี่ยนสินค้า   & อยากเปลี่ยนไซซ์จาก M เป็น L ได้ไหมคะ \\
\hline
Delivery Delay          & ส่งช้า            & สั่งไปตั้งแต่อาทิตย์ที่แล้วแต่ของยังไม่มาเลยครับ \\
\hline
Address Change          & เปลี่ยนที่อยู่    & ขอเปลี่ยนที่อยู่จัดส่งเป็นที่ทำงานแทนได้ไหมคะ \\
\hline
Refund                  & ขอคืนเงิน         & อยากขอเงินคืนเพราะยกเลิกออเดอร์แล้วคะ \\
\hline
Excess Item             & ได้รับสินค้าเกิน  & ได้รับสินค้าเกินจากที่สั่งไป 1 ชิ้นครับ \\
\hline
Shipping Status Inquiry & สอบถามสถานะขนส่ง  & ติดตามสถานะขนส่งยังไงคะ / ออเดอร์จะมาถึงวันไหนคะ / ติดต่อคนขับส่งของไม่ได้เลยคะ \\

\end{longtable}

\vspace{0.4cm}


\subsection{จริยธรรมข้อมูล (Data Ethics)}

เนื่องจากโครงงานนี้ใช้ข้อมูลจริงจากลูกค้าร้านค้าออนไลน์ จึงกำหนดแนวทางปฏิบัติด้านจริยธรรมข้อมูลอย่างเคร่งครัด ดังนี้:

\subsubsection{การลบข้อมูลระบุตัวตน (Anonymization/De-identification)}
ข้อมูลส่วนบุคคล (PII) เช่น ชื่อ เบอร์โทรศัพท์ ที่อยู่ และหมายเลขคำสั่งซื้อ จะถูกลบหรือเข้ารหัสก่อนนำเข้าสู่กระบวนการวิเคราะห์ เพื่อให้สอดคล้องกับ พ.ร.บ. คุ้มครองข้อมูลส่วนบุคคล (PDPA)

\subsubsection{การเก็บรักษาและสิทธิการเข้าถึง}
ข้อมูลจะถูกจัดเก็บในสภาพแวดล้อมที่ปลอดภัย (Secure Environment) จำกัดสิทธิ์การเข้าถึงเฉพาะผู้ดำเนินงานวิจัย และจะถูกลบทิ้งอย่างถาวรเมื่อเสร็จสิ้นโครงการ

\subsubsection{ข้อตกลงการใช้งาน}
ข้อมูลถูกนำมาใช้เพื่อวัตถุประสงค์ในการศึกษาวิจัยและพัฒนาระบบอัตโนมัติภายในเท่านั้น จะไม่มีการนำข้อมูลดิบไปแสวงหาผลกำไรหรือเผยแพร่สู่สาธารณะโดยเด็ดขาด

\subsection{โมเดลเปรียบเทียบ}

\subsubsection{โมเดล Large Language Model (Zero-shot)}
\begin{enumerate}[label=(\arabic*)]
    \item Typhoon2-7B: เลือกใช้เนื่องจากเป็นโมเดลที่พัฒนาและฝึกสอนสำหรับภาษาไทยโดยเฉพาะ (Thai-centric) จึงคาดว่าจะมีความเข้าใจบริบทภาษาพูดและคำสแลงไทยได้ดี
    \item GPT-4o: เลือกใช้เพื่อเป็นตัวแทนของโมเดลเชิงพาณิชย์ระดับ State-of-the-Art ที่มีประสิทธิภาพสูงในการวิเคราะห์และจัดหมวดหมู่ข้อความ
    \item Claude 3.5 Sonnet: เลือกใช้เนื่องจากมีความโดดเด่นด้านการปฏิบัติตามคำสั่ง (Instruction Following) ซึ่งเหมาะสมกับการใช้ Prompt ที่มีความซับซ้อน
\end{enumerate}

\subsubsection{โมเดล Supervised Baseline}
\begin{enumerate}[label=(\arabic*)]
    \item WangchanBERTa: เลือกใช้เป็นตัวแทนของโมเดล Pre-trained ภาษาไทยแบบดั้งเดิมที่เหมาะสำหรับงาน text classification
    \item mBERT: เลือกใช้เป็นตัวแทนของโมเดล Multilingual เพื่อเปรียบเทียบประสิทธิภาพกับโมเดลภาษาไทยเฉพาะทาง
\end{enumerate}

\subsection{ระเบียบวิธีทดสอบและตัวชี้วัด (Testing Protocol and Metrics)}

\subsubsection{ระเบียบวิธีทดสอบ (Testing Protocol)}
ข้อมูลจะถูกแบ่งด้วยเทคนิค Stratified Split เป็นสัดส่วน Train 70\%, Validation 15\%, และ Test 15\% เพื่อรักษาการกระจายตัวของข้อมูลที่ไม่สมดุล (Imbalanced data) ในทุกชุด การประเมินประสิทธิภาพจะวัดจากชุด Test set โดยใช้ WangchanBERTa เป็น Baseline สำหรับเปรียบเทียบกับ LLM ในการประเมินฝั่ง Supervised Model จะมีการรันแบบสุ่มซ้ำอย่างน้อย 3 รอบ และรายงานผลเป็นค่าเฉลี่ยเพื่อความน่าเชื่อถือเชิงสถิติ

\subsubsection{ตัวชี้วัด (Metrics)}
\begin{enumerate}[label=(\arabic*)]
    \item Accuracy: วัดความแม่นยำรวมของการจำแนกประเภทข้อความทั้งหมด
    \item Macro F1-score: ใช้เป็นตัวชี้วัดหลักเพื่อรับมือกับปัญหาคลาสไม่สมดุล เนื่องจากพิจารณาความสำคัญของทุกประเภทข้อความเท่าเทียมกัน
    \item Per-class F1-score: เพื่อวิเคราะห์ประสิทธิภาพเชิงลึกว่าโมเดลทำได้ดีหรือด้อยในประเภทคำร้องเรียนใดเป็นพิเศษ
    \item Confusion Matrix: ใช้สรุปและแสดงจุดที่โมเดลมักทำนายผิดพลาดข้ามคลาส (Misclassification) ซึ่งช่วยให้เข้าใจความกำกวมของคำร้องเรียน
\end{enumerate}

% -------------------------------------------------------
\section{สถานที่ทำวิจัย}
% -------------------------------------------------------

สาขาวิชาวิทยาการคอมพิวเตอร์ วิทยาลัยการคอมพิวเตอร์ มหาวิทยาลัยขอนแก่น หรืออื่น ๆ

% -------------------------------------------------------
\section{ผลที่คาดว่าจะได้รับ}
% -------------------------------------------------------

\begin{enumerate}
    \item ชุดข้อมูลข้อความติดต่อของลูกค้าที่มีการติด label พร้อมสำหรับใช้งานและต่อยอด
    \item โมเดลต้นแบบสำหรับจำแนกประเภทข้อความร้องเรียนของลูกค้าโดยใช้ Zero-shot Text Classification
    \item ผลการเปรียบเทียบประสิทธิภาพระหว่างโมเดล LLM
\end{enumerate}

% -------------------------------------------------------
\section{แผนการดำเนินงาน}
% -------------------------------------------------------

\noindent\textbf{ตารางที่~2} แผนการดำเนินงาน

\vspace{0.3cm}
\noindent
{%
\renewcommand{\arraystretch}{1.5}%
\setlength{\tabcolsep}{3pt}%
\small
\setlength{\LTleft}{0pt}
\begin{longtable}{|>{\centering\arraybackslash}p{3.8cm}|*{10}{>{\centering\arraybackslash}p{0.72cm}|}}
\hline
\multicolumn{1}{|>{\centering\arraybackslash}p{3.8cm}|}{\textbf{การดำเนินงาน}} &
\multicolumn{10}{c|}{\textbf{ปี 2569}} \\
\cline{2-11}
 & \textbf{1} & \textbf{2} & \textbf{3} & \textbf{4} & \textbf{5} & \textbf{6} & \textbf{7} & \textbf{8} & \textbf{9} & \textbf{10} \\
\hline
\endfirsthead
\multicolumn{11}{@{}l@{}}{\textbf{ตารางที่~2} แผนการดำเนินงาน (ต่อ)} \\[2pt]
\hline
\multicolumn{1}{|>{\centering\arraybackslash}p{3.8cm}|}{\textbf{การดำเนินงาน}} &
\multicolumn{10}{c|}{\textbf{ปี 2569}} \\
\cline{2-11}
 & \textbf{1} & \textbf{2} & \textbf{3} & \textbf{4} & \textbf{5} & \textbf{6} & \textbf{7} & \textbf{8} & \textbf{9} & \textbf{10} \\
\hline
\endhead
\hline
\endfoot
ศึกษางานวิจัยที่เกี่ยวข้อง
 & \cellcolor[RGB]{210,120,80} & \cellcolor[RGB]{210,120,80} & & & & & & & & \\
\hline
เก็บรวบรวมและเตรียมข้อมูล
 & & & \cellcolor[RGB]{210,120,80} & \cellcolor[RGB]{210,120,80} & & & & & & \\
\hline
ออกแบบและพัฒนาระบบ
 & & & & & \cellcolor[RGB]{210,120,80} & \cellcolor[RGB]{210,120,80} & & & & \\
\hline
{\small ทดสอบด้วย Zero-shot\newline Classification ผ่าน LLM}
 & & & & & & & \cellcolor[RGB]{210,120,80} & \cellcolor[RGB]{210,120,80} & & \\
\hline
ประเมินผลและเปรียบเทียบโมเดล
 & & & & & & & & & \cellcolor[RGB]{210,120,80} & \\
\hline
{\small สรุปผลและจัดทำ\newline รายงานฉบับสมบูรณ์}
 & & & & & & & & & & \cellcolor[RGB]{210,120,80} \\
\hline
\end{longtable}%
}


% -------------------------------------------------------
\section{งบประมาณ}
% -------------------------------------------------------

\noindent หมวดวัสดุอุปกรณ์

\begin{enumerate}[label=(\arabic*), leftmargin=2\firstindentlength, topsep=0pt, itemsep=0pt]
    \item ค่าวัสดุสำนักงาน (กระดาษ ปากกา ฯลฯ)
    \item ค่าวัสดุคอมพิวเตอร์/สื่อบันทึก (แฟลชไดรฟ์ ฯลฯ)
\end{enumerate}

\noindent หมวดค่าใช้สอย

\begin{enumerate}[label=(\arabic*), leftmargin=2\firstindentlength, topsep=0pt, itemsep=0pt]
    \item ค่าถ่ายเอกสาร
    \item ค่าจัดรูปเล่ม
\end{enumerate}

\noindent หมวดค่าใช้จ่ายอื่น ๆ

\begin{enumerate}[label=(\arabic*), leftmargin=2\firstindentlength, topsep=0pt, itemsep=0pt]
    \item ค่าเดินทาง/ประสานงานเก็บข้อมูล
    \item ค่าตอบแทนผู้ทดสอบ (User Study ขนาดเล็ก)
\end{enumerate}
